\section{相关工作}\label{sec:related}
\textbf{联合调度—内存优化。}COSMA~\citep{li2023combined} 已使用 ILP 联合优化算子调度、内存分配与张量替换，故联合优化本身属于既有工作，本文不主张“首次联合”。本文的场景更窄也更细粒度：多缓存微操作 DAG、评测器定义的静态有备份/生成溢出类别、依赖前沿生产启发式，以及将启发式与机器验证最优值相连接的固定序连续地址证书。\textbf{寄存器压力与重物化。}指令调度与寄存器分配的联合优化~\citep{goodman1988integrated,shobaki2013preallocation,shobaki2022register} 同样在容量约束下塑造活跃区间；Checkmate~\citep{jain2019checkmate} 离线求解重物化调度，DTR~\citep{kirisame2020dtr} 是其面向动态模型的在线对应。二者均以重计算换取内存，而本文基准禁止重计算，要求为静态核函数 DAG 生成显式的溢出/重载操作与具体地址——这正是证书必须包含打包验证的原因。\textbf{缓存与替换。}Belady 的固定请求流规则~\citep{belady1966study} 与考虑对象大小、代价的缓存研究~\citep{coester2022blockaware} 已处理非一致代价，干净感知替换在编译器内存管理中亦有先例~\citep{liberatore1999evaluation,guo2004belady}，故本文不主张一般加权缓存意义上的新颖性。本文的贡献在于其具体组合：优化器自身可改变的拓扑请求序、多个片上缓存、标准的一次/两次传输溢出代价，以及“下界—物理打包”桥梁。无溢出编译~\citep{chatarasi2026spillfree} 与本文互补：本文的规划器针对评测器仍要求溢出流量的实例。
